\par\Needspace{6\baselineskip}\medskip
\begin{center}\hypertarget{native-198-article148}{}\pdfbookmark[2]{148.}{article-148}148.\end{center}
Quoties elementa e tribus observationibus derivanda adhuc penitus incognita
sunt (cui casui methodus nostra imprimis accomodata est), in hypothesi
prima ut iam monuimus pro $P$ et $Q$ valores approximati $\dfrac{\theta''}{\theta}$ et $\theta\theta''$
accipientur, ubi $\theta$ et $\theta''$ aliquantisper ex intervallis temporum non correctis
derivandae sunt. Quorum ratione ad interva\'{l}la correcta per $\mu:1$ et $\mu'':1$
resp. expressa, habebimus in hypothesi prima
\[
\begin{aligned}
X&=\log\mu-\log\mu''+\log\eta-\log\eta''\\
Y&=-\log\mu-\log\mu''-\log\eta-\log\eta''\\
&\quad+\mathrm{Comp.}\log\cos f+\mathrm{Comp.}\log\cos f'\\
&\quad+\mathrm{Comp.}\log\cos f''+2\log r'-\log r-\log r''.
\end{aligned}
\]
Logarithmi quantitatum $\mu,\mu''$ respectu partium reliquarum nullius sunt momenti;
$\log\eta$ et $\log\eta''$, qui ambo sunt positivi, in $X$ aliquatenus se invicem
destruunt, praesertim quoties temporum intervalla fere aequalia sunt, unde
$X$ valorem exiguum modo positivum modo negativum obtinet; contra in $Y$ e
partibus negativis $\log\eta$ et $\log\eta''$ compensatio quidem aliqua partium positivarum
$\mathrm{Comp.}\log\cos f$, $\mathrm{Comp.}\log\cos f'$, $\mathrm{Comp.}\log\cos f''$ oritur, sed minus perfecta,
plerumque enim hae illas notabiliter superant. De signo ipsius $\log\dfrac{r'r'}{rr''}$ in genere
nihil determinare licet.

Iam quoties motus heliocentricus inter observationes modicus est, raro
opus erit, usque ad hypothesin quartam progredi: plerumque tertia, saepius
\SourceLeaf{199}iam secunda praecisionem sufficientem praestabit, quin adeo interdum numeris
ex ipsa hypothesi prima resultantibus acquiescere licebit. Iuvabit semper, ad
maiorem minoremve praecisionis gradum, qua observationes gaudent, respicere:
ingratum enim foret opus, in calculo praecisionem affectare centies milliesve
maiorem ea quam observationes permittunt. In his vero rebus iudicium per
exercitationem frequentem practicam melius quam per praecepta acuitur, peritique
facile acquirent facultatem quandam ubi consistere conveniat recte diiudicandi.
